\section{Tradeoffs and Engineering Decisions}

\subsection{Option Comparison}
\begin{table}[H]
    \centering
    \small
    \begin{tabular}{p{2.5cm} c c c p{3.5cm}}
        \toprule
        \textbf{Option} & \textbf{Perf.} & \textbf{Risk} & \textbf{Cost} & \textbf{Reasoning} \\
        \midrule
        Option A & High & Medium & High & Best technical margin but more integration effort. \\
        Option B & Medium & Low & Medium & Balanced path for cost and schedule. \\
        Option C & Medium & High & Low & Attractive cost profile but poor robustness margin. \\
        \bottomrule
    \end{tabular}
    \caption{Representative architecture tradeoff table}
    \label{tab:tradeoffs}
\end{table}

\subsection{Sizing or Budget Snapshot}
Use this section for first-order engineering calculations that justify the chosen
direction:

\begin{align}
    P_{budget} &= V_{in} \times I_{avg} \\
    Margin_{thermal} &= T_{limit} - T_{measured} \\
    Headroom_{timing} &= t_{allowed} - t_{observed}
\end{align}

\highlightbox{
Keep calculations close to the design decision they support. Detached
mathematics is easy to ignore during review.
}

\subsection{Key Decisions}
\begin{itemize}
    \item identify the selected concept and why it won;
    \item state what evidence would invalidate the current choice;
    \item capture which design parameter is most sensitive to requirement changes.
\end{itemize}
